% Section 11.4, from lectures/L11/appendix/c_exercises.md. The intro's "görs lämpligen på passet"
% (exercise 9 needs the board) reads "görs lämpligen i samband med Labb 3".

\section{Övningar}
\label{c11:sec:exercises}\label{c11:app:c}

\begin{aside}
\textbf{Så kontrollerar du ditt arbete.} Övningarna märkta \textbf{Program} skrivs i Microchip
Studio och testas först i simulatorn och sedan, om ni har kortet, på ett Arduino Uno.
Övning~\ref{c11:ex:9}, kontrollen, kräver kortet och en multimeter, och görs lämpligen i samband
med Labb~3.

Övningarna märkta \textbf{Räkna för hand}, \textbf{Förståelse} och \textbf{Konstruktion} görs på
papper, utan miniräknare, och kontrolleras mot bilaga~\ref{app:answers}. Gör varje övning innan du
läser lösningen.

Varje övning är märkt med sin sort. Exakt en övning är en \textbf{Kontroll}, och i det här kapitlet
är det övning~\ref{c11:ex:9}.
\end{aside}

\exercise{Ingångar, beslut och utgångar}{Förståelse}
\label{c11:ex:1}
Beskriv för vart och ett av styrsystemen nedan vilka \textbf{ingångar} det har, vilka
\textbf{utgångar}, och vilket \textbf{beslut} det fattar, som i \secref{c11:a1}.
\begin{parts}
  \item En garageport med en knapp, en motor som kan köra upp och ned, och två
    gränslägesbrytare, en som känner att porten är helt öppen och en att den är helt stängd.
  \item En handtork som blåser när en IR-givare ser en hand, och slutar tio sekunder efter att
    handen försvunnit.
  \item Övergångsstället i slutuppgiften, \hyperref[lab3:b]{del~B} i Labb~3.
\end{parts}

\exercise{Förkopplingsmotståndet}{Räkna för hand}
\label{c11:ex:2}
Stiftet ger 5~V. Räkna med 2~V över en röd lysdiod och 3~V över en grön.
\begin{parts}
  \item Hur stor blir strömmen genom en röd lysdiod med 220~Ω, 330~Ω och 1~kΩ?
  \item Hur stor blir strömmen genom en grön lysdiod med 220~Ω? Varför lyser den kanske svagare än
    den röda med samma motstånd?
  \item Vilket motstånd behövs för 10~mA genom en röd lysdiod?
  \item Någon tar ett motstånd på 47~Ω. Hur stor blir strömmen, och vad säger databladets gräns om
    det?
\end{parts}

\exercise{Hela porten}{Räkna för hand}
\label{c11:ex:3}
Sex röda lysdioder sitter på PB0--PB5, var och en med 220~Ω, och alla lyser samtidigt.
\begin{parts}
  \item Hur stor är den sammanlagda strömmen ut ur porten?
  \item Jämför med gränsen på ungefär 100~mA för en grupp av stift, från \secref{c11:a2}. Hur mycket
    marginal finns det?
  \item Varför är 330~Ω ett säkrare val än 220~Ω om många lysdioder kan lysa samtidigt?
\end{parts}

\exercise{Reläet}{Förståelse}
\label{c11:ex:4}
Titta på kopplingen i \secref{c11:a3}.
\begin{parts}
  \item Varför kopplas reläspolen inte direkt till stiftet?
  \item Vad gör dioden, och vad händer om den saknas?
  \item Mikrodatorn och 24~V-kretsen har ingen gemensam ledning. Varför är det en fördel?
  \item Någon kopplar lampan på 24~V direkt mellan stiftet och jord, utan relä. Vad händer, med
    lampan och med kortet?
\end{parts}

\exercise{Knappar som studsar}{Förståelse}
\label{c11:ex:5}
\begin{parts}
  \item Förklara med egna ord varför programmet i \secref{c11:a4} räknar 6 i stället för 2 när
    väntan tas bort.
  \item Varför väntar programmet 20~ms, och inte 1~ms eller 500~ms? Vad går fel med var och en av
    de två andra?
  \item Lysdioden i \code{led_follows_button.asm} i kapitel~\ref{c10:ch} (\secref{c10:b4}) tänds så
    länge knappen hålls ned, och har ingen avstudsning alls. Varför behövs den inte där?
\end{parts}

\exercise{Hur lång tid tar ett varv?}{Räkna för hand}
\label{c11:ex:6}
Använd formeln för \code{delay_s} i \secref{c11:b4}.
\begin{parts}
  \item Hur många cykler tar ett anrop av \code{delay_s} med \code{r24} = 4?
  \item Vägarbetsljusets loop har, utöver de två anropen, tre instruktioner per tillstånd
    (\code{ldi}, \code{out}, \code{ldi}) och ett \code{rjmp}. Hur många cykler tar ett helt varv?
    Jämför med de 128\,000\,716 cykler som kursens kontroll mätte.
  \item Hur mycket för långt blir ett varv, i mikrosekunder? Hur mycket blir det på en timme?
  \item Jämför med resonatorns noggrannhet i \secref{c11:a7}. Vilken av de två felkällorna är värd
    att bry sig om?
\end{parts}

\exercise{Nattläge}{Konstruktion}
\label{c11:ex:7}
På natten stängs trafikljus ofta av, och bara den gula lampan blinkar: en halv sekund tänd, en
halv sekund släckt.
\begin{parts}
  \item Skriv tillståndstabellen, som i \secref{c11:b2}.
  \item Rita flödesplanen.
  \item Skriv programmet, med den gula lysdioden på PB1 och subrutinen \code{delay_ms}. Kan
    \code{r24} laddas en enda gång, före loopen? Varför, eller varför inte?
  \item Hur lång blir en hel blinkperiod, räknat i cykler?
\end{parts}

\exercise{En knapp som växlar}{Program}
\label{c11:ex:8}
Skriv ett program där varje tryck på knappen på PD2 \textbf{växlar} lysdioden på PB5: ett tryck
tänder, nästa släcker, och så vidare. Använd avstudsning som i \secref{c11:a4}.
\begin{parts}
  \item Rita flödesplanen.
  \item Skriv programmet och testa det i simulatorn.
  \item Vad händer på kortet om väntan efter trycket tas bort? Prova, om ni har kortet.
\end{parts}

\exercise{Kontroll: strömmen genom lysdioden}{Kontroll}
\label{c11:ex:9}
\emph{Räkna för hand, mät på kortet, förklara skillnaden.}

En röd lysdiod med 220~Ω sitter på stift 8 (PB0), och ett program tänder den.

\task{a) För hand}
Räkna ut strömmen genom lysdioden, med 2~V över den.

\task{b)}
\textbf{Mät spänningen över motståndet} med en multimeter, när lysdioden lyser. Räkna ut strömmen
med Ohms lag. Mät också spänningen mellan stiftet och jord.

\task{c) Jämför}
Stämmer strömmarna i a) och b)? Förklara skillnaden med det du mätte: är spänningen på stiftet
verkligen 5~V, och är spänningen över lysdioden verkligen 2~V?

\task{d)}
Byt till en grön lysdiod och gör om b). Vad ändras, och varför?

\exercise{Felsökning på kortet}{Förståelse}
\label{c11:ex:10}
Programmet för trafikljuset fungerar i simulatorn. Förklara för vart och ett av symptomen nedan
vad det troligaste felet är, och hur du kontrollerar det.
\begin{parts}
  \item Den gröna lampan lyser aldrig. De andra två fungerar.
  \item Alla tre lamporna ser ut att lysa svagt, hela tiden.
  \item Ingenting händer alls, och utskriften efter överföringen slutar med
    \code{ser_open(): can't open device}.
  \item Den röda och den gula lampan har bytt plats: gul lyser när röd ska lysa.
\end{parts}

\exercise{En dag på ett trafikljus}{Räkna för hand, Fördjupning}
\label{c11:ex:11}
\begin{parts}
  \item Kortets resonator går 0,5~\% för fort. Hur mycket fel har en klocka byggd på programmets
    fördröjningar efter ett dygn?
  \item Hur många hela ljuscykler, om 8~s var, hinner trafikljuset i slutuppgiften med på ett
    dygn? Räkna med exakt 8~s.
\end{parts}
